\oldpage[IV-2]{84}

un corps de caractéristique $p>0$, $k$ le corps des fractions
rationnelles $k_0(s,t)$ à $2$ indéterminées, $L$ le corps
$k(s^{1/p},t^{1/p})$ et $X=\operatorname{Spec}(L)$. Soient d'autre part $u$
une troisième indéterminée et soit $K$ le corps
$k(u,s^{1/p}+ut^{1/p})$; on vérifie aisément que $k$ est algébriquement fermé
dans $K$ et que $L\otimes_kK$ a un nilradical non nul, donc $X_{(K)}$ n'est
pas réduit; mais comme $(X_{(K)})_{\mathrm{red}}$ n'est pas défini sur $k$,
et que $K$ ne contient aucune sous-extension de degré fini sur $k$ autre que
$k$, le plus petit corps de définition de $(X_{(K)})_{\mathrm{red}}$ ne peut
être fini sur $k$.

\item[(ii)] Étant donné un élément de l'un des ensembles d'arrivée des
applications (4.8.2), et une extension $K_1$ de $K$, il est équivalent de dire
que l'objet considéré admet un plus petit corps de définition ou que son image
par $\varepsilon_{K_1,K}$ admet un plus petit sous-corps de définition (les
deux sous-corps étant nécessairement les mêmes), comme on l'a vu dans
(4.8.4).
\end{enumerate}
\end{remarks}
\subsection{Corps de définition d'une partie d'un préschéma}

\begin{proposition}[4.9.1]
\label{IV.2.4.9.1-fr}
Soient $k$ un corps, $X$ un $k$-préschéma, $K$ une extension de $k$,
$\overline T$ une partie fermée de $X_{(K)}$, $K'$ une sous-extension de $K$.
Les conditions suivantes sont équivalentes:
\begin{enumerate}
\item[a)] $\overline T$ est définie sur $K'$.
\item[b)] La partie ouverte $\overline U=X_{(K)}-\overline T$ de $X_{(K)}$
est définie sur $K'$.
\item[b')] Le sous-préschéma de $X_{(K)}$ induit sur l'ouvert $\overline U$
est défini sur $K'$.
\item[c)] Il existe un sous-préschéma fermé $\overline Z$ de $X_{(K)}$ ayant
$\overline T$ pour espace sous-jacent, et défini sur $K'$.
\end{enumerate}
\end{proposition}

\begin{proof}
Si $g:X_{(K)}\to X_{(K')}$ est la projection canonique, dire que
$\overline T$ est définie sur $K'$ signifie qu'il existe une partie fermée
$T'$ de $X_{(K')}$ telle que $\overline T=g^{-1}(T')$; comme
$g^{-1}(X_{(K')}-T')=X_{(K)}-g^{-1}(T')$ pour toute partie $T'$ de
$X_{(K')}$, a), b) et b') sont équivalentes. Si $\overline Z$, ayant
$\overline T$ pour espace sous-jacent, est défini sur $K'$, on a
$\overline Z=Z'\otimes_{K'}K$, où $Z'$ est un sous-préschéma fermé de
$X_{(K')}$, et si $T'$ est le sous-espace sous-jacent à $Z'$, on sait que
$\overline T=g^{-1}(T')$ (I, 4.4.1), donc $\overline T$ est définie sur
$K'$. Inversement, pour voir que a) entraîne c), il suffit de considérer un
sous-préschéma fermé $Z'$ de $X_{(K')}$ ayant pour espace sous-jacent $T'$
(I, 5.2.1) et de prendre $\overline Z=Z'\otimes_{K'}K$.
\end{proof}

\begin{corollary}[4.9.2]
\label{IV.2.4.9.2-fr}
Avec les notations de (4.9.1), supposons que $K'$ soit un corps de définition
de $\overline T$; toute extension $K''\subset K'$ de $k$, telle que $K'$ soit
extension radicielle de $K''$, est aussi un corps de définition de
$\overline T$.
\end{corollary}

\begin{proof}
Il suffit d'observer que la projection canonique
$g:X_{(K')}\to X_{(K'')}$ est un morphisme entier, surjectif et radiciel,
donc un homéomorphisme (2.4.5).
\end{proof}

\begin{remark}[4.9.3]
\label{IV.2.4.9.3-fr}
Il résulte de (4.9.2) qu'une partie fermée $\overline T$ de $X_{(K)}$ n'admet
pas nécessairement de \emph{plus petit} corps de définition. Par exemple, si
$K=k(t)$, où $t$ est une indéterminée et $k$ un corps parfait de
caractéristique $p>0$, une partie fermée $\overline F$ de $X_{(K)}$ ne peut
avoir de plus petit corps de définition que si elle est déjà définie sur $k$;
en effet, pour tout sous-corps $K'$ de $K$ contenant $k$ et $\ne k$, on a
$K'^p\ne K'$, $K'^p$ contient $k$ et $K'$
\oldpage[IV-2]{85}
est une extension radicielle de $K'^p$, donc si $K'$ est corps de définition
de $\overline F$, il en est de même de $K'^p$. Toutefois:
\end{remark}

\begin{proposition}[4.9.4]
\label{IV.2.4.9.4-fr}
Avec les notations de \textup{(4.9.1)}, supposons que $K$ soit une extension séparable
de $K'$; pour que $K'$ soit un corps de définition de $\overline T$, il faut
et il suffit que $K'$ soit un corps de définition du sous-préschéma réduit de
$X_{(K)}$ ayant $\overline T$ pour espace sous-jacent.
\end{proposition}

\begin{proof}
La condition est évidemment suffisante en vertu de (4.9.1); pour voir qu'elle
est nécessaire, notons que si $\overline T=g^{-1}(T')$, où $T'$ est une
partie fermée de $X_{(K')}$, et si $Z'$ est le sous-préschéma réduit de
$X_{(K')}$ ayant $T'$ pour espace sous-jacent, alors
$\overline Z=Z'\otimes_{K'}K=\operatorname{Spec}(K)\times_{K'}Z'$ est réduit
en vertu du critère (4.6.1, \emph{e)}) et de l'hypothèse sur $K$ et a pour espace
sous-jacent $T$.
\end{proof}

\begin{corollary}[4.9.5]
\label{IV.2.4.9.5-fr}
Supposons vérifiée l'une des hypothèses suivantes:
\begin{enumerate}
\item[\textup{a)}] $k$ est de caractéristique $0$.
\item[\textup{b)}] $K$ est une extension algébrique de $k$ et $X$ est de type fini sur
$k$.
\end{enumerate}
Alors toute partie fermée $\overline T$ de $X_{(K)}$ possède un plus petit
corps de définition, qui est une extension séparable de $k$ (et finie sur $k$
dans l'hypothèse \textup{b))}.
\end{corollary}

\begin{proof}
Dans l'hypothèse \emph{a)}, $K$ est séparable sur chacun de ses sous-corps, et il
résulte donc de (4.9.4) que les corps de définition de $\overline T$ sont les
mêmes que ceux du sous-préschéma réduit de $X_{(K)}$ ayant $\overline T$ pour
espace sous-jacent. Le corollaire résulte alors de (4.8.11). Dans l'hypothèse
\emph{b)} on sait que $K$ est une extension radicielle d'une extension séparable
$K_1$ de $k$. Alors, pour toute sous-extension $K'$ de $K$, $K'$ est
radicielle sur $K'\cap K_1$, donc, en vertu de (4.9.2), $K'$ est un corps de
définition de $\overline T$ si et seulement si $K'\cap K_1$ l'est. Cela nous
ramène au cas où $K'=K_1$, donc au cas où $K'$ est une extension \emph{algébrique
séparable} de $k$, et on achève le raisonnement comme dans le cas \emph{a)}.
\end{proof}

\begin{proposition}[4.9.6]
\label{IV.2.4.9.6-fr}
Soient $k$ un corps, $X$ un $k$-préschéma, $K$ une extension de $k$,
$\overline U$ une partie à la fois ouverte et fermée de $X_{(K)}$,
$\overline V$ son complémentaire. Soit $K'$ une sous-extension de $K$;
alors les conditions suivantes sont équivalentes:
\begin{enumerate}
\item[\textup{a)}] $K'$ est un corps de définition de la partie $\overline U$ de
$X_{(K)}$.
\item[\textup{a')}] $K'$ est un corps de définition du sous-préschéma fermé de
$X_{(K)}$ induit sur l'ouvert $\overline U$.
\item[\textup{b)}] $K'$ est un corps de définition de la partie $\overline V$ de
$X_{(K)}$.
\item[\textup{b')}] $K'$ est un corps de définition du sous-préschéma fermé de
$X_{(K)}$ induit sur l'ouvert $\overline V$.
\end{enumerate}
\end{proposition}

\begin{proof}
Il est clair que \emph{a')} entraîne \emph{a)}, et \emph{a)} entraîne \emph{b')} en vertu de (4.9.1).
Comme $\overline U$ et $\overline V$ jouent des rôles symétriques, \emph{b')}
entraîne \emph{b)} et \emph{b)} entraîne \emph{a')}, ce qui achève la démonstration.
\end{proof}

\begin{corollary}[4.9.7]
\label{IV.2.4.9.7-fr}
Sous les hypothèses de \textup{(4.9.6)}, il existe un plus petit corps de définition
$K'$ pour la partie $\overline U$ de $X_{(K)}$, et $K'$ est aussi le plus
petit corps de définition de la partie $\overline V$ de $X_{(K)}$ et des
sous-préschémas fermés de $X_{(K)}$ induits sur les ouverts $\overline U$ et
$\overline V$. Si en outre $X$ est de type fini sur $k$, $K'$ est une
extension finie séparable de $k$.
\end{corollary}

\begin{proof}
Compte tenu de (4.9.6) et (4.8.11), il n'y a à démontrer que la dernière
assertion. En vertu de (4.8.15, (ii)), on peut se borner au cas où $K$ est
algébriquement clos. Soit alors $\overline k$ la clôture algébrique de $k$
contenue dans $K$, et désignons par $W_i$ les composantes
\oldpage[IV-2]{86}
connexes de $X_{(\overline k)}$ ($1\leq i\leq m$); si
$p:X_{(K)}\to X_{(\overline k)}$ est la projection canonique, on sait que les
$p^{-1}(W_i)$ sont les composantes connexes de $X_{(K)}$ (4.4.6), donc
$U$ est une réunion d'un certain nombre de ces composantes, et est
par suite défini sur $\overline k$; la conclusion résulte par suite de
(4.9.5, \emph{b)}).
\end{proof}
\section{Dimension, profondeur, régularité dans les préschémas localement
noethériens}

Ce paragraphe se borne à reprendre, dans le langage géométrique, et avec
divers compléments de nature technique, les notions et résultats d'Algèbre
commutative exposés dans le chapitre \textbf{0}, \S\S~16 et 17.

\subsection{Dimension des préschémas.}

\begin{env}[5.1.1]
\label{IV.2.5.1.1-fr}
Soient $A$ un anneau, $\mathfrak J$ un idéal de $A$. Rappelons
$(\mathbf{0},16.1)$ que les définitions de la théorie de la dimension d'un anneau
$(\mathbf{0},16.1.1)$ et celles de la théorie de la dimension combinatoire des espaces
topologiques $(\mathbf{0},14.1.2)$ donnent les relations suivantes:
\begin{align}
\dim(\operatorname{Spec}(A))&=\dim(A)
\tag{5.1.1.1}\label{IV.2.5.1.1.1-fr}\\
\dim(V(\mathfrak J))&=\dim(A/\mathfrak J)
\tag{5.1.1.2}\label{IV.2.5.1.1.2-fr}\\
\operatorname{codim}(V(\mathfrak J),\operatorname{Spec}(A))
&=\operatorname{ht}(\mathfrak J)
\tag{5.1.1.3}\label{IV.2.5.1.1.3-fr}
\end{align}
\par\noindent\textbf{(5.1.1.4)} $\operatorname{Spec}(A)$ est un espace
caténaire $\Longleftrightarrow A$ est un anneau caténaire.
\par\noindent\textbf{(5.1.1.5)} $\operatorname{Spec}(A)$ est
équidimensionnel $\Longleftrightarrow A$ est équidimensionnel
$\Longleftrightarrow$ les idéaux premiers minimaux $\mathfrak p_\alpha$ de
$A$ sont tels que les dimensions des anneaux $A/\mathfrak p_\alpha$ soient
toutes égales.
\par\noindent\textbf{(5.1.1.6)} $\operatorname{Spec}(A)$ est
équicodimensionnel $\Longleftrightarrow A$ est équicodimensionnel
$\Longleftrightarrow$ tous les idéaux maximaux de $A$ ont même hauteur.

On rappelle qu'un anneau noethérien $A$ est dit \emph{biéquidimensionnel} si
$\operatorname{Spec}(A)$ est biéquidimensionnel, c'est-à-dire si
$\operatorname{Spec}(A)$ est noethérien et $A$ est équidimensionnel,
équicodimensionnel, caténaire et de dimension finie.
\end{env}

\begin{proposition}[5.1.2]
\label{IV.2.5.1.2-fr}
Soient $X$ un préschéma, $Y$ une partie fermée irréductible de $X$, $y$ le
point générique de $Y$. On a
\begin{equation}
\operatorname{codim}(Y,X)=\dim(\mathcal O_{X,y})
\tag{5.1.2.1}\label{IV.2.5.1.2.1-fr}
\end{equation}
\end{proposition}

\begin{proof}
En effet (\textbf{I}, 2.4.2) les parties fermées irréductibles de $X$ contenant $y$
sont canoniquement en correspondance biunivoque avec les parties fermées
irréductibles de $\operatorname{Spec}(\mathcal O_{X,y})$, donc sont en
correspondance biunivoque avec les idéaux premiers de $\mathcal O_{X,y}$, et
(5.1.2.1) résulte des définitions.

En particulier, on retrouve ainsi le fait que les points génériques des
composantes irréductibles de $X$ sont les \emph{seuls} points $x\in X$ tels
que $\dim(\mathcal O_x)=0$ (\textbf{I}, 1.1.14).
\end{proof}

\oldpage[IV-2]{87}

\begin{corollary}[5.1.3]
\label{IV.2.5.1.3-fr}
Pour toute partie fermée $Y$ d'un préschéma $X$, on a
\begin{equation}
\operatorname{codim}(Y,X)=\inf_{y\in Y}\dim(\mathcal O_{X,y}).
\tag{5.1.3.1}\label{IV.2.5.1.3.1-fr}
\end{equation}
En outre, si $X$ est localement noethérien, on a, pour tout $x\in Y$
\begin{equation}
\operatorname{codim}_x(Y,X)=
\inf_{\substack{y\in Y,\,x\in\overline{\{y\}}}}
\dim(\mathcal O_{X,y}).
\tag{5.1.3.2}\label{IV.2.5.1.3.2-fr}
\end{equation}
\end{corollary}

\begin{proof}
La relation (5.1.3.2) résulte en effet de (5.1.3.1) et de $(\mathbf{0},14.2.6)$.

Ce corollaire permet de \emph{définir}, pour une partie \emph{quelconque}
$Y$ d'un préschéma $X$, la \emph{codimension} $\operatorname{codim}(Y,X)$ de
$Y$ dans $X$ comme égale au second membre de (5.1.3.1).
\end{proof}

\begin{proposition}[5.1.4]
\label{IV.2.5.1.4-fr}
Pour tout préschéma $X$, on a
\begin{equation}
\dim(X)=\sup_{x\in X}(\dim(\mathcal O_x)).
\tag{5.1.4.1}\label{IV.2.5.1.4.1-fr}
\end{equation}
Si toute partie fermée irréductible de $X$ contient un point fermé, on a aussi
\begin{equation}
\dim(X)=\sup_{x\in F}(\dim(\mathcal O_x))
\tag{5.1.4.2}\label{IV.2.5.1.4.2-fr}
\end{equation}
où $F$ est l'ensemble des points fermés de $X$.
\end{proposition}

\begin{proof}
Il suffit de remarquer (en vertu de (\textbf{I}, 2.4.2)) que les chaînes de parties
fermées irréductibles de $X$ correspondent biunivoquement aux chaînes de
parties fermées irréductibles d'un schéma local de $X$; lorsque toute partie
fermée irréductible de $X$ contient un point fermé, on peut évidemment se
borner aux schémas locaux aux points fermés de $X$.

On notera que toute partie fermée irréductible de $X$ contient un point fermé
lorsque $X$ est \emph{quasi-compact} $(\mathbf{0}_{\mathrm I},2.1.3)$; nous verrons un
peu plus loin (5.1.11) qu'il en est de même lorsque $X$ est \emph{localement
noethérien}.
\end{proof}

\begin{corollary}[5.1.5]
\label{IV.2.5.1.5-fr}
Pour qu'un préschéma $X$ soit caténaire, il faut et il suffit que, pour tout
$x\in X$, l'anneau local $\mathcal O_x$ soit caténaire. Si de plus toute
partie fermée irréductible de $X$ contient un point fermé, il suffit, pour que
$X$ soit caténaire, que $\mathcal O_x$ soit caténaire pour tout point fermé
$x$ de $X$.
\end{corollary}

\begin{proof}
Le raisonnement est le même que dans (5.1.4), puisqu'il s'agit de comparer des
chaînes de parties fermées irréductibles ayant mêmes extrémités.
\end{proof}

\begin{env}[5.1.6]
\label{IV.2.5.1.6-fr}
Nous allons maintenant examiner les propriétés plus spéciales liées à des
hypothèses noethériennes.

Rappelons $(\mathbf{0},16.2.3)$ qu'un anneau \emph{local noethérien} $A\ne0$ est de
dimension \emph{finie}, égale au nombre minimum de générateurs d'un idéal de
définition de $A$; pour tout idéal premier $\mathfrak p$ de $A$, la hauteur de
$\mathfrak p$, égale à $\dim(A_{\mathfrak p})$, est donc aussi finie. Ces
propriétés, jointes à (5.1.2), (5.1.3) et (5.1.4), montrent que:
\end{env}

\begin{proposition}[5.1.7]
\label{IV.2.5.1.7-fr}
Pour toute partie fermée non vide $Y$ d'un préschéma localement noethérien
$X$, $\operatorname{codim}(Y,X)$ est finie. Si $X$ est noethérien et affine
et $Y$ une partie fermée irréductible de $X$, $\operatorname{codim}(Y,X)$ est
égal au nombre minimum des sections $s_i$ de $\mathcal O_X$ au-dessus de $X$
telles que $Y$ soit une composante irréductible de l'ensemble des $x\in X$
tels que $s_i(x)=0$ pour tout $i$.
\end{proposition}

\begin{corollary}[5.1.8]
\label{IV.2.5.1.8-fr}
Soient $X$ un préschéma localement noethérien, $\mathcal L$ un
$\mathcal O_X$-Module inversible, $f$ une section de $\mathcal L$ au-dessus
de $X$. Alors toute composante irréductible de l'ensemble $Z$
\oldpage[IV-2]{88}
des $x\in X$ tels que $f(x)=0$ est de codimension $\leq1$ dans $X$; elle est
de codimension $1$ si $Z$ ne contient aucune composante irréductible de $X$.
\end{corollary}

\begin{proof}
On peut se borner au cas où $\mathcal L=\mathcal O_X$. Si $y$ est un point
générique d'une composante irréductible $Y$ de $Z$, l'idéal $(f_y)$ de
$\mathcal O_{X,y}$ doit être tel que $\mathcal O_{X,y}/(f_y)$ n'ait qu'un
seul idéal premier, ce qui signifie que $f_y$ engendre un idéal de définition
de l'anneau local noethérien $\mathcal O_{X,y}$; on a donc
$\operatorname{codim}(Y,X)\leq1$ (5.1.7); si $Z$ ne contient aucune
composante irréductible de $X$, on ne peut avoir
$\operatorname{codim}(Y,X)=0$ en vertu de $(\mathbf{0},14.2.1)$.
\end{proof}

\begin{proposition}[5.1.9]
\label{IV.2.5.1.9-fr}
Soient $X$ un préschéma localement noethérien, $Y$ une partie fermée de $X$.
Si $x\in Y$ est tel que $\mathcal O_{X,x}$ soit un anneau caténaire, on a
\begin{align}
\operatorname{codim}_x(Y,X)
&=\dim(\mathcal O_{X,x})\setmin\operatorname{codim}(\overline{\{x\}},Y)\notag\\
&=\dim(\mathcal O_{X,x})\setmin\dim(\mathcal O_{Y,x})
\tag{5.1.9.1}\label{IV.2.5.1.9.1-fr}
\end{align}
\end{proposition}

\begin{proof}
Soient $Y_i$ ($1\leq i\leq n$) les composantes irréductibles de $Y$
contenant $x$ (qui sont en nombre \emph{fini} puisque $Y$ est localement
noethérien), et soit $y_i$ le point générique de $Y_i$. Si l'on pose
$A=\mathcal O_{X,x}$, et si $\mathfrak p_i$ est l'idéal premier de $A$
correspondant à $Y_i\cap\operatorname{Spec}(A)$, les parties fermées
irréductibles de $Y$ contenant $x$ correspondent biunivoquement aux idéaux
premiers de $A$ contenant un des $\mathfrak p_i$, donc
$\dim(\mathcal O_{Y,x})=\sup_i\dim(A/\mathfrak p_i)$; d'autre part, on a
$\mathcal O_{X,y_i}=A_{\mathfrak p_i}$, et l'hypothèse sur $A$ entraîne
$\dim(A)=\dim(A/\mathfrak p_i)+\dim(A_{\mathfrak p_i})$ $(\mathbf{0},16.1.4)$; la
conclusion résulte donc de la relation
$\operatorname{codim}_x(Y,X)=\inf_i(\dim(\mathcal O_{X,y_i}))$ ((5.1.2) et
$(\mathbf{0},14.2.6)$).
\end{proof}

\begin{proposition}[5.1.10]
\label{IV.2.5.1.10-fr}
\begin{enumerate}
\item[\textup{(i)}] Dans tout préschéma $X$, toute partie localement constructible non
vide $E$ contient un point $x$ tel que $\{x\}$ soit une partie localement
fermée de $X$ \textup{(ou, ce qui revient au même, tel que $x$ soit \emph{isolé dans}
$\overline{\{x\}}$)}.
\item[\textup{(ii)}] Soient $X$ un préschéma localement noethérien, $x$ un point de
$X$ tel que $\{x\}$ soit localement fermé dans $X$; alors on a
$\dim(\overline{\{x\}})\leq1$; tout point $y\ne x$ de
$\overline{\{x\}}$ est par suite fermé dans $X$.
\end{enumerate}
\end{proposition}

\begin{proof}
(i) Le résultat est un cas particulier du lemme suivant:
\end{proof}

\begin{lemma}[5.1.10.1]
\label{IV.2.5.1.10.1-fr}
Soit $X$ un espace topologique ayant la propriété suivante: pour toute partie
localement fermée $Z\ne\varnothing$ de $X$, il existe une partie $Z'$ de $Z$,
localement fermée dans $X$ \textup{(ou \emph{dans} $Z$, ce qui revient au même)} et un point
$x\in Z'$, fermé \textup{dans} $Z'$. Alors toute partie localement fermée
$Z\ne\varnothing$ de $X$ contient un point $x$ isolé dans
$\overline{\{x\}}$.
\end{lemma}

\begin{proof}
En effet, soit $Z'\subset Z$ une partie localement fermée de $Z$ contenant un
point $x$ tel que $x$ soit fermé \emph{dans} $Z'$. Il y a un voisinage ouvert
$U$ de $x$ dans $X$ tel que $Z'\cap U$ soit fermé \emph{dans} $U$, donc $x$ est aussi
fermé dans $U$; cela signifie que $U\cap\overline{\{x\}}=\{x\}$, autrement
dit que $x$ est isolé dans $\overline{\{x\}}$.

Le lemme s'applique à tout espace sous-jacent à un préschéma $X$, car $Z$ est
alors aussi espace sous-jacent à un préschéma (\textbf{I}, 5.2.1) et il suffit de
prendre pour $Z'$ un ouvert affine \emph{dans} $Z$, qui est un espace de Kolmogoroff
quasi-compact, donc contient un point fermé $(\mathbf{0}_{\mathrm I},2.1.3)$.
\end{proof}

\oldpage[IV-2]{89}

\begin{proof}
\textup{(ii)} Soit $Z$ le sous-préschéma réduit de $X$ ayant
$\overline{\{x\}}$ pour espace
sous-jacent; l'hypothèse entraîne que $\{x\}$ est ouvert \emph{dans} $Z$;
donc, pour tout $z\in Z$, le point générique $x$ de
$\operatorname{Spec}(\mathcal O_{Z,z})$ est \emph{isolé} dans
$\operatorname{Spec}(\mathcal O_{Z,z})$; mais l'anneau
$A=\mathcal O_{Z,z}$ est un anneau local noethérien intègre, et l'hypothèse
entraîne qu'il existe un $f\in A$ tel que $A_f$ soit un \emph{corps}; on sait
$(\mathbf{0},16.3.3)$ que cela entraîne $\dim(A)\leq1$. Comme
$\dim(\mathcal O_{Z,z})\leq1$ pour tout $z\in Z$, on a bien $\dim(Z)\leq1$.
\end{proof}

\begin{corollary}[5.1.11]
\label{IV.2.5.1.11-fr}
Si $X$ est un préschéma localement noethérien, toute partie fermée non vide de
$X$ contient un point fermé.
\end{corollary}

\begin{proof}
En effet, toute partie fermée de $X$ est constructible
$(\mathbf{0}_{\mathrm{III}},9.1.1\text{ et }9.1.5)$ et il suffit d'appliquer (5.1.10).
\end{proof}

\begin{env}[5.1.12]
\label{IV.2.5.1.12-fr}
Soient $X$ un préschéma, $\mathcal F$ un $\mathcal O_X$-Module
quasi-cohérent \emph{de type fini}, $S=\operatorname{Supp}(\mathcal F)$ son
support, qui est \emph{fermé} dans $X$ $(\mathbf{0}_{\mathrm I},5.2.2)$. Si, pour tout
$x\in X$, on considère $\operatorname{Supp}(\mathcal F_x)$ comme une partie
fermée du schéma local $\operatorname{Spec}(\mathcal O_x)$, on a, par
définition $(\mathbf{0},16.1.7)$
$\dim(\mathcal F_x)=\dim(\operatorname{Supp}(\mathcal F_x))$; mais on a
\[
\operatorname{Supp}(\mathcal F_x)
=S\cap\operatorname{Spec}(\mathcal O_{X,x})
=\operatorname{Spec}(\mathcal O_{S,x})
\]
où, dans le dernier terme, $S$ désigne un sous-préschéma fermé quelconque de
$X$ ayant $S$ pour espace sous-jacent. Si l'on remarque que les composantes
irréductibles de $\operatorname{Spec}(\mathcal O_{S,x})$ sont les
intersections de $\operatorname{Spec}(\mathcal O_{X,x})$ et des composantes
irréductibles de $S$ contenant $x$, et correspondent biunivoquement à ces
dernières, on voit que
\begin{equation}
\dim(\mathcal F_x)=\dim(\mathcal O_{S,x})
\tag{5.1.12.1}\label{IV.2.5.1.12.1-fr}
\end{equation}
en vertu de (5.1.1.1); il résulte aussi de (5.1.2) que l'on a
\begin{equation}
\dim(\mathcal F_x)=\operatorname{codim}(\overline{\{x\}},S)
\tag{5.1.12.2}\label{IV.2.5.1.12.2-fr}
\end{equation}
si $X$ est localement noethérien.

On dit que $\mathcal F$ est \emph{équidimensionnel au point} $x\in X$ si
$\mathcal F_x$ est un $\mathcal O_{X,x}$-module équidimensionnel,
c'est-à-dire $(\mathbf{0},16.1.7)$ si $\operatorname{Supp}(\mathcal F_x)$ est
équidimensionnel en tant que partie fermée de
$\operatorname{Spec}(\mathcal O_{X,x})$; il revient au même de dire que
l'anneau $\mathcal O_{S,x}$ est \emph{équidimensionnel}.

On appelle \emph{dimension de} $\mathcal F$ et l'on note $\dim(\mathcal F)$
la dimension du \emph{support} $\operatorname{Supp}(\mathcal F)$; il résulte
de (5.1.4) et (5.1.12.1) que l'on a
\begin{equation}
\dim(\mathcal F)=\sup_{x\in X}\dim(\mathcal F_x).
\tag{5.1.12.3}\label{IV.2.5.1.12.3-fr}
\end{equation}
Si $X=\operatorname{Spec}(A)$ est un schéma affine et si
$\mathcal F=\widetilde M$, où $M$ est un $A$-module de type fini, on a
$\dim(\mathcal F)=\dim(M)$ en vertu de $(\mathbf{0},16.1.7)$ et (5.1.4).
\end{env}

\begin{proposition}[5.1.13]
\label{IV.2.5.1.13-fr}
Soient $X$ un préschéma, $\mathcal F$ un $\mathcal O_X$-Module
quasi-cohérent de type fini, $x$ un point de $X$, $x'$ une générisation de
$x$ dans $X$; on a alors
\begin{equation}
\dim(\mathcal F_{x'})\leq\dim(\mathcal F_x)
\tag{5.1.13.1}\label{IV.2.5.1.13.1-fr}
\end{equation}
\end{proposition}

\begin{proof}
Cela résulte aussitôt de (5.1.12.1) et des définitions.
\end{proof}

\oldpage[IV-2]{90}
\subsection{Dimension d'un préschéma algébrique.}

\begin{proposition}[5.2.1]
\label{IV.2.5.2.1-fr}
Soient $k$ un corps, $X$ un préschéma irréductible localement de type fini sur
$k$, $\xi$ son point générique. Alors $X$ est biéquidimensionnel et l'on a
$\dim(X)=\operatorname{deg.tr}_k k(\xi)$.
\end{proposition}

\begin{proof}
Tout anneau local $\mathcal O_x$ est anneau local d'une $k$-algèbre de type
fini, donc quotient d'un anneau local régulier $(0,17.3.9)$, et l'on sait
$(0,16.5.12)$ que ces anneaux sont caténaires; par suite $X$ est un espace
caténaire (5.1.5) et comme $X$ est irréductible, il suffit (5.1.1) de montrer
que pour \emph{tout point fermé} $x\in X$, on a
\begin{equation}
\dim(\mathcal O_x)=\operatorname{deg.tr}_k k(\xi).
\tag{5.2.1.1}\label{IV.2.5.2.1.1-fr}
\end{equation}
On peut évidemment supposer $X$ réduit et affine, donc \emph{intègre}
d'anneau $A$, algèbre de type fini sur $k$. Soit
$n=\operatorname{deg.tr}_k k(\xi)$, $k(\xi)=K$ étant le corps des fractions
de $A$. On sait (Bourbaki, \emph{Alg. comm.}, chap.~V, \S~3,
n\textsuperscript{o}~1, th.~1) qu'il existe une sous-$k$-algèbre
$B=k[t_1,\ldots,t_n]$ de $A$, où les $t_i$ sont algébriquement indépendants
sur $k$, telle que $A$ soit une $B$-algèbre \emph{finie}. Soit
$\mathfrak m=\mathfrak j_x$, qui est par hypothèse un idéal maximal de $A$;
$\mathfrak n=B\cap\mathfrak m$ est donc un idéal maximal de $B$ (Bourbaki,
\emph{Alg. comm.}, chap.~V, \S~2, n\textsuperscript{o}~1, prop.~1), et
$A_{\mathfrak m}$ est un anneau local de la $B_{\mathfrak n}$-algèbre
\emph{finie} $S^{-1}A$, où $S=B-\mathfrak n$; comme $B_{\mathfrak n}$ est
intégralement clos et $S^{-1}A$ intègre, on a
$\dim(A_{\mathfrak m})=\dim(B_{\mathfrak n})$ $(0,16.1.6)$. On peut donc se
borner au cas où $A=k[t_1,\ldots,t_n]$; on sait alors que
$k'=k(x)$ est une extension \emph{finie} de $k$ (I, 6.4.2); soit
$A'=k'[t_1,\ldots,t_n]$; compte tenu de (I, 2.4.6 et 3.3.14), il existe un
idéal maximal $\mathfrak m'$ de $A'$ au-dessus de $\mathfrak m$ et tel que
$k'$ soit le corps résiduel de $A'_{\mathfrak m'}$. Comme $A'$ est intègre et
est une $A$-algèbre finie, le même raisonnement que ci-dessus montre que
$\dim(A'_{\mathfrak m'})=\dim(A_{\mathfrak m})$, si bien qu'on est ramené au
cas où $A/\mathfrak m=k$. Or dans ce cas $\mathfrak m$ est engendré par des
polynômes $t_i-a_i$ (où $a_i\in k$, $1\leq i\leq n$, les $a_i$ étant les
images canoniques des $t_i$ dans $A/\mathfrak m$). Remplaçant $t_i$ par
$t_i-a_i$, on voit finalement qu'on est ramené au cas où
$\mathfrak m=\sum_{i=1}^nAt_i$. Le complété de l'anneau local
$A_{\mathfrak m}$ est alors l'anneau de séries formelles
$k[[t_1,\ldots,t_n]]$; on sait qu'il a même dimension que $A_{\mathfrak m}$
$(0,16.2.4)$, et d'autre part la dimension de $k[[t_1,\ldots,t_n]]$ est égale
à $n$ $(0,17.1.4)$; d'où la conclusion.
\end{proof}

\begin{corollary}[5.2.2]
\label{IV.2.5.2.2-fr}
Pour un préschéma $X$ localement de type fini sur un corps $k$, $\dim(X)$
coïncide avec le nombre défini dans (4.1.1).
\end{corollary}

\begin{proof}
En effet, si $X_\alpha$ sont les sous-préschémas réduits ayant pour espaces
sous-jacents les composantes irréductibles de $X$, on a
$\dim(X)=\sup_\alpha(\dim(X_\alpha))$ $(0,14.1.2.1)$ et il suffit d'appliquer
(5.2.1) aux $X_\alpha$.
\end{proof}

\begin{corollary}[5.2.3]
\EGAReaderTarget{ega4.c04.cor.5-2-3}{IV.2.5.2.3-fr}
Soient $X$ un préschéma localement de type fini sur un corps $k$, $x$ un
point de $X$. On a
\begin{equation}
\dim_x(X)=\dim(\mathcal O_x)+\operatorname{deg.tr}_k k(x).
\tag{5.2.3.1}\label{IV.2.5.2.3.1-fr}
\end{equation}
\end{corollary}

\begin{proof}
On sait qu'il existe un voisinage ouvert $U$ de $x$ dans $X$ tel que
$\dim_x(X)=\dim(U)$ $(0,14.1.4.1)$, et l'on peut supposer que les composantes
irréductibles de $U$ soient les $X_i\cap U$, où les $X_i$ sont les
composantes irréductibles de $X$ contenant $x$; comme $U\cap X_i$

